% GENERATED FILE. Edit canonical lessons, then run npm run generate:books.
% canonical-source-hash: fnv1a64:b8c99f6733530aa1
% canonical-lessons: BN-C512-unun, BN-C512-khunti, BN-C512-sarashi, BN-C512-dhakna, BN-C512-bata

\chapter{Kitchen Tools}
\label{ch:bn-kitchen-tools}
\hlchaptermodality{512}

\begin{chapteropening}
\textbf{By the end of this chapter:} \emph{I can name some kitchen tools and say what they do.}

Complete the last lesson of chapter 512: I can name some kitchen tools and say what they do.
\end{chapteropening}

\section[\texorpdfstring{\bn{উনুন} (unun)}{উনুন (unun)}]{\bn{উনুন} (unun) — a stove (of clay); an oven}

\label{lesson:BN-C512-unun}



\begin{quote}
Before the new one: say the Bengali for to pour, then the Bengali for to dip.
\end{quote}

\subsection*{You'll want to know: \bn{উনুন}}
\textbf{\bn{উনুন}} — \emph{unun} — \textquotedblleft{}a stove (of clay); an oven\textquotedblright{}.

\begin{quote}
\bn{দিদিমা} \bn{উনুনে} \bn{ভাত} \bn{রান্না} \bn{করেন।}
\end{quote}

\begin{quote}
\emph{Didimā unune bhāt rānnā kôren.}
\end{quote}

\begin{quote}
My grandmother cooks rice on the clay stove.
\end{quote}

\begin{cousinweb}
\textbf{\bn{উনুন}} (\emph{unun}) is a tadbhava: CDIAL lists Bengali \emph{unun}, \textquotedblleft{}oven, fireplace\textquotedblright{}, under Sanskrit \emph{uṣṇayati}, \textquotedblleft{}heats\textquotedblright{} (CDIAL 2394).
\end{cousinweb}

\begin{grammarlens}[title={What you've built}]
One more word for describing what you lived and felt.
\end{grammarlens}

\subsection*{Your turn}
\begin{itemize}
\raggedright
  \item \emph{Say it:} \emph{unun}
  \item \emph{Say it:} \emph{unun}, once more
  \item \emph{Say it:} \emph{Didimā unune bhāt rānnā kôren.} — My grandmother cooks rice on the clay stove.
  \item \emph{Recall:} say the Bengali for to pour, then the Bengali for to dip, then say \emph{unun} again
\end{itemize}

\begin{tcolorbox}[breakable,colback=teal!4,colframe=teal!35!black,title={Before you move on}]
What does \bn{উনুন} mean? (\textquotedblleft{}A stove (of clay); an oven\textquotedblright{}.) Say it once more.
\end{tcolorbox}
\section[\texorpdfstring{\bn{খুন্তি} (khunti)}{খুন্তি (khunti)}]{\bn{খুন্তি} (khunti) — a spatula; a flat cooking spoon}

\label{lesson:BN-C512-khunti}



\begin{quote}
Before the new one: say the Bengali for to dip, then the Bengali for a stove.
\end{quote}

\subsection*{You'll want to know: \bn{খুন্তি}}
\textbf{\bn{খুন্তি}} — \emph{khunti} — \textquotedblleft{}a spatula; a flat cooking spoon\textquotedblright{}.

\begin{quote}
\bn{আমি} \bn{খুন্তিটা} \bn{কড়াইয়ের} \bn{পাশে} \bn{রাখলাম।}
\end{quote}

\begin{quote}
\emph{Āmi khuntiṭā kôṛāiyer pāshe rākhlām.}
\end{quote}

\begin{quote}
I put the spatula beside the wok.
\end{quote}

\begin{cousinweb}
\textbf{\bn{খুন্তি}} (\emph{khunti}) is a tadbhava from Sanskrit \emph{khanitra}, \textquotedblleft{}a digging tool\textquotedblright{}; CDIAL lists Bengali \emph{khunti} for a long-handled digging tool (CDIAL 3814). The kitchen sense, a spatula, is the modern language's own.
\end{cousinweb}

\begin{grammarlens}[title={What you've built}]
One more word for describing what you lived and felt.
\end{grammarlens}

\subsection*{Your turn}
\begin{itemize}
\raggedright
  \item \emph{Say it:} \emph{khunti}
  \item \emph{Say it:} \emph{khunti}, once more
  \item \emph{Say it:} \emph{Āmi khuntiṭā kôṛāiyer pāshe rākhlām.} — I put the spatula beside the wok.
  \item \emph{Recall:} say the Bengali for to dip, then the Bengali for a stove, then say \emph{khunti} again
\end{itemize}

\begin{tcolorbox}[breakable,colback=teal!4,colframe=teal!35!black,title={Before you move on}]
What does \bn{খুন্তি} mean? (\textquotedblleft{}A spatula; a flat cooking spoon\textquotedblright{}.) Say it once more.
\end{tcolorbox}
\section[\texorpdfstring{\bn{সাঁড়াশি} (sā̃ṛāshi)}{সাঁড়াশি (sā̃ṛāshi)}]{\bn{সাঁড়াশি} (sā̃ṛāshi) — tongs}

\label{lesson:BN-C512-sarashi}



\begin{quote}
Before the new one: say the Bengali for a stove, then the Bengali for a spatula.
\end{quote}

\subsection*{You'll want to know: \bn{সাঁড়াশি}}
\textbf{\bn{সাঁড়াশি}} — \emph{sā̃ṛāshi} — \textquotedblleft{}tongs\textquotedblright{}.

\begin{quote}
\bn{আমি} \bn{সাঁড়াশিতে} \bn{ধরে} \bn{গরম} \bn{হাঁড়িটা} \bn{নামালাম।}
\end{quote}

\begin{quote}
\emph{Āmi sā̃ṛāshite dhôre gôrôm hā̃ṛiṭā nāmālām.}
\end{quote}

\begin{quote}
I gripped the hot pot with the tongs and took it off the stove.
\end{quote}

\begin{cousinweb}
\textbf{\bn{সাঁড়াশি}} (\emph{sā̃ṛāshi}) is a tadbhava: CDIAL lists Bengali \emph{sā̃ṛāsi}, \textquotedblleft{}pincers\textquotedblright{}, under Sanskrit \emph{saṁdaṁśa}, \textquotedblleft{}tongs\textquotedblright{} (CDIAL 12897).
\end{cousinweb}

\begin{grammarlens}[title={What you've built}]
One more word for describing what you lived and felt.
\end{grammarlens}

\subsection*{Your turn}
\begin{itemize}
\raggedright
  \item \emph{Say it:} \emph{sā̃ṛāshi}
  \item \emph{Say it:} \emph{sā̃ṛāshi}, once more
  \item \emph{Say it:} \emph{Āmi sā̃ṛāshite dhôre gôrôm hā̃ṛiṭā nāmālām.} — I gripped the hot pot with the tongs and took it off the stove.
  \item \emph{Recall:} say the Bengali for a stove, then the Bengali for a spatula, then say \emph{sā̃ṛāshi} again
\end{itemize}

\begin{tcolorbox}[breakable,colback=teal!4,colframe=teal!35!black,title={Before you move on}]
What does \bn{সাঁড়াশি} mean? (\textquotedblleft{}Tongs\textquotedblright{}.) Say it once more.
\end{tcolorbox}
\section[\texorpdfstring{\bn{ঢাকনা} (ḍhāknā)}{ঢাকনা (ḍhāknā)}]{\bn{ঢাকনা} (ḍhāknā) — a lid}

\label{lesson:BN-C512-dhakna}



\begin{quote}
Before the new one: say the Bengali for a spatula, then the Bengali for tongs.
\end{quote}

\subsection*{You'll want to know: \bn{ঢাকনা}}
\textbf{\bn{ঢাকনা}} — \emph{ḍhāknā} — \textquotedblleft{}a lid\textquotedblright{}.

\begin{quote}
\bn{হাঁড়ির} \bn{ঢাকনাটা} \bn{কোথায়}?
\end{quote}

\begin{quote}
\emph{Hā̃ṛir ḍhāknāṭā kothāy?}
\end{quote}

\begin{quote}
Where is the lid of the pot?
\end{quote}

\begin{cousinweb}
\textbf{\bn{ঢাকনা}} (\emph{ḍhāknā}) comes from a reconstructed root \emph{ḍhakk}, \textquotedblleft{}cover\textquotedblright{}; CDIAL lists Bengali \emph{ḍhākan}, \textquotedblleft{}lid\textquotedblright{} (CDIAL 5574).
\end{cousinweb}

\begin{grammarlens}[title={What you've built}]
One more word for describing what you lived and felt.
\end{grammarlens}

\subsection*{Your turn}
\begin{itemize}
\raggedright
  \item \emph{Say it:} \emph{ḍhāknā}
  \item \emph{Say it:} \emph{ḍhāknā}, once more
  \item \emph{Say it:} \emph{Hā̃ṛir ḍhāknāṭā kothāy?} — Where is the lid of the pot?
  \item \emph{Recall:} say the Bengali for a spatula, then the Bengali for tongs, then say \emph{ḍhāknā} again
\end{itemize}

\begin{tcolorbox}[breakable,colback=teal!4,colframe=teal!35!black,title={Before you move on}]
What does \bn{ঢাকনা} mean? (\textquotedblleft{}A lid\textquotedblright{}.) Say it once more.
\end{tcolorbox}
\section[\texorpdfstring{\bn{বাটা} (bāṭā)}{বাটা (bāṭā)}]{\bn{বাটা} (bāṭā) — to grind (spices) on a stone}

\label{lesson:BN-C512-bata}



\begin{quote}
Before the new one: say the Bengali for tongs, then the Bengali for a lid.
\end{quote}

\subsection*{You'll want to know: \bn{বাটা}}
\textbf{\bn{বাটা}} — \emph{bāṭā} — \textquotedblleft{}to grind (spices) on a stone\textquotedblright{}.

\begin{quote}
\bn{দিদিমা} \bn{শিলনোড়ায়} \bn{মশলা} \bn{বাটলেন।}
\end{quote}

\begin{quote}
\emph{Didimā shilnoṛāy môshlā bāṭlen.}
\end{quote}

\begin{quote}
My grandmother ground the spices on the grinding stone.
\end{quote}

Past (he/she …): \bn{তিনি} \bn{আদা} \bn{বাটলেন}, \emph{tini ādā bāṭlen}.

\begin{cousinweb}
\textbf{\bn{বাটা}} (\emph{bāṭā}) is a tadbhava: CDIAL lists Bengali \emph{bāṭā}, \textquotedblleft{}to pound, to crush\textquotedblright{}, under Sanskrit \emph{vartayati}, \textquotedblleft{}makes turn\textquotedblright{} (CDIAL 11356).
\end{cousinweb}

\begin{grammarlens}[title={What you've built}]
One more word for describing what you lived and felt.
\end{grammarlens}

\subsection*{Your turn}
\begin{itemize}
\raggedright
  \item \emph{Say it:} \emph{bāṭā}
  \item \emph{Say it:} \emph{bāṭā}, once more
  \item \emph{Say it:} \emph{Didimā shilnoṛāy môshlā bāṭlen.} — My grandmother ground the spices on the grinding stone.
  \item \emph{Recall:} say the Bengali for tongs, then the Bengali for a lid, then say \emph{bāṭā} again
\end{itemize}

\begin{tcolorbox}[breakable,colback=teal!4,colframe=teal!35!black,title={Before you move on}]
What does \bn{বাটা} mean? (\textquotedblleft{}To grind (spices) on a stone\textquotedblright{}.) Say it once more.
\end{tcolorbox}
